\chapter{实验过程与结果}

\section{实验配置}

本实验使用仓库中的示例策略和示例明文进行验证。指数抽象路线和真实 pairing 路线均使用 \verb|Project| 目录下的同一组示例文件。两条路线采用相同访问策略和明文内容，便于对照底层群运算替换前后的行为。明文内容为：

\begin{lstlisting}[caption={示例明文}]
Hello LSSS-ABE.
\end{lstlisting}

图 \ref{fig:screenshot-message-body} 展示了示例明文文件的实际内容。后续成功解密场景会将输出文件与该内容进行对照。

\begin{figure}[H]
\centering
\includegraphics[width=0.72\textwidth]{examle_message.png}
\caption{原始明文 message.txt 截图}
\label{fig:screenshot-message-body}
\end{figure}

访问策略为：

\[
    (A\wedge B)\vee C.
\]

实验设计遵循两个原则。第一，测试场景必须覆盖访问结构的不同语义，包括 AND 分支、OR 分支、一般阈值门和不满足策略的情况。第二，测试不能只验证正常路径，还要覆盖工程边界，例如特殊属性名、密文篡改、重复属性和非法策略输入。这样可以为系统并非只对示例策略“刚好可用”提供实验依据，并展示其对常见边界输入的处理能力。

这些场景的选择与参考资料中的关键要求一一对应：Bethencourt 等人~[2]强调 CP-ABE 的访问策略由属性集合满足；Waters~[3]要求根据 $I=\{i:\rho(i)\in S\}$ 求重构常数；Liu 等人~[4]强调阈值访问树不应只局限于 AND/OR；实验三课件~[1]要求检验属性基加密算法实现是否正确。因此，测试既要有满足策略的成功样例，也必须有不满足策略、阈值门、重复属性和非法输入等反例或边界样例。

测试属性集合与边界检查设置如表 \ref{tab:test-scenarios} 所示。基础指数版包含九个场景；真实 pairing 版复用其中的核心语义，并在 BLS12-381 群元素上重新验证成功解密、拒绝解密、特殊属性名、密文篡改、一般阈值门、重复属性和非法策略输入。

\begin{table}[H]
\centering
\caption{测试场景设置}
\label{tab:test-scenarios}
\begin{tabular}{>{\raggedright\arraybackslash}p{0.10\textwidth} >{\raggedright\arraybackslash}p{0.25\textwidth} >{\raggedright\arraybackslash}p{0.29\textwidth} >{\raggedright\arraybackslash}p{0.20\textwidth}}
\toprule
场景 & 输入或操作 & 策略或对象 & 实验现象 \\
\midrule
1 & $\{A,B\}$ & $(A\wedge B)\vee C$ & 解密成功 \\
2 & $\{C\}$ & $(A\wedge B)\vee C$ & 解密成功 \\
3 & $\{A\}$ & $(A\wedge B)\vee C$ & 解密失败 \\
4 & \verb|role:admin| & 单属性策略 & 安全生成文件并解密成功 \\
5 & 篡改 \verb|payload.tag| & 密文 JSON & HMAC 拒绝 \\
6 & $\{D,F\}$ & $2$-of-$3(D,E,F)$ & 解密成功 \\
7 & $\{D\}$ & $2$-of-$3(D,E,F)$ & 解密失败 \\
8 & $\{X\}$ & 重复属性策略 $X\wedge X$ & 解密成功 \\
9 & \verb|k=0|、\verb|k>n|、缺失字段、空属性等 & 非法策略输入 & 加密前拒绝 \\
\bottomrule
\end{tabular}
\end{table}

\section{实验流程}

完整实验流程包括以下环节：

\begin{enumerate}
    \item 运行 Setup，生成公共参数和主密钥。
    \item 分别为 $\{A,B\}$、$\{C\}$、$\{A\}$ 生成用户私钥。
    \item 使用公共参数和访问策略加密示例明文。
    \item 使用 $\{A,B\}$ 和 $\{C\}$ 私钥解密，并与原始明文进行比对。
    \item 使用 $\{A\}$ 私钥尝试解密，记录不满足策略时的拒绝行为。
    \item 使用 \verb|role:admin| 验证属性名包含冒号时仍能生成安全私钥文件并解密。
    \item 篡改密文中的 \verb|payload.tag|，验证 HMAC 完整性检查会拒绝密文。
    \item 使用 $2$-of-$3$ 阈值策略验证一般阈值门的成功和失败路径。
    \item 使用重复属性策略验证多行映射到同一属性时的重构行为。
    \item 使用非法策略验证输入校验能在加密前拒绝错误访问树。
\end{enumerate}

上述流程在实验中由烟雾测试脚本统一执行。基础指数版脚本使用 \path{tests/artifacts} 目录重新生成密钥、密文和输出文件；真实 pairing 版脚本使用独立的 \path{tests/artifacts_pairing} 目录保存测试产物。这样可以保证两条路线的测试结果均来自各自完整流程，也避免指数版产物与 pairing 版产物混用。相比零散命令输出，完整的 smoke 输出更能反映项目整体复现情况。

烟雾测试脚本命令如下：

\begin{lstlisting}[language=bash,caption={烟雾测试命令}]
Set-Location .\Project
powershell -NoProfile -ExecutionPolicy Bypass -File .\tests\smoke.ps1
\end{lstlisting}

真实 pairing 路线使用独立脚本验证。该脚本调用同一个 PowerShell CLI，但增加 \verb|-Adapter| \verb|pairing| 参数，使命令转发到 Python BLS12-381 后端：

\begin{lstlisting}[language=bash,caption={真实 pairing 版烟雾测试命令}]
Set-Location .\Project
powershell -NoProfile -ExecutionPolicy Bypass -File .\tests\smoke_pairing.ps1
\end{lstlisting}

基础指数版关键输出如下：

\begin{lstlisting}[caption={基础指数版烟雾测试关键输出}]
Scenario 1 OK: {A,B} decrypted successfully.
Scenario 2 OK: {C} decrypted successfully.
Scenario 3 OK: {A} correctly rejected.
Scenario 4 OK: filename-safe attribute key for role:admin decrypted successfully.
Scenario 5 OK: tampered ciphertext correctly rejected.
Scenario 6 OK: {D,F} satisfied 2-of-3 threshold policy.
Scenario 7 OK: {D} correctly rejected by 2-of-3 threshold policy.
Scenario 8 OK: repeated attribute rows decrypted successfully.
Scenario 9 OK: invalid policies correctly rejected.
All smoke tests passed.
\end{lstlisting}

真实 pairing 版关键输出如下：

\begin{lstlisting}[caption={真实 pairing 版烟雾测试关键输出}]
Pairing scenario 1 OK: {A,B} decrypted successfully.
Pairing scenario 2 OK: {C} decrypted successfully.
Pairing scenario 3 OK: {A} correctly rejected.
Pairing scenario 4 OK: role:admin decrypted successfully.
Pairing scenario 5 OK: tampered ciphertext correctly rejected.
Pairing scenario 6 OK: 2-of-3 threshold behavior is correct.
Pairing scenario 7 OK: repeated attribute rows decrypted successfully.
Pairing scenario 8 OK: invalid policies correctly rejected.
All pairing smoke tests passed.
\end{lstlisting}

图 \ref{fig:screenshot-smoke-body} 展示了基础指数版真实终端中的烟雾测试完整输出，截图包含九个 OK 场景和最终通过提示。图 \ref{fig:screenshot-pairing-smoke-body} 展示了真实 pairing 路线的完整终端输出，截图中包含 conda 环境激活、进入项目源码目录、八类 pairing 场景和最终通过提示。两张截图分别对应指数抽象路线和 BLS12-381 真实配对路线，说明同一访问结构在两种底层群运算实现下均完成了成功解密、拒绝解密和边界输入验证。

\begin{figure}[H]
\centering
\includegraphics[width=0.92\textwidth]{All_Smoke_Tests_Pass.png}
\caption{基础指数版烟雾测试完整输出截图}
\label{fig:screenshot-smoke-body}
\end{figure}

\begin{figure}[H]
\centering
\includegraphics[width=0.95\textwidth]{pairing_smoke_tests_pass.png}
\caption{真实 pairing 路线烟雾测试完整输出截图}
\label{fig:screenshot-pairing-smoke-body}
\end{figure}

除统一烟雾测试外，本实验也保留了逐条运行 Setup、KeyGen、Encrypt 和 Decrypt 命令的结果。图 \ref{fig:screenshot-setup-body} 至图 \ref{fig:screenshot-decrypt-body} 展示了四个核心命令的实际运行截图。

\begin{figure}[H]
\centering
\includegraphics[width=0.95\textwidth]{setup.png}
\caption{Setup 命令运行截图}
\label{fig:screenshot-setup-body}
\end{figure}

\begin{figure}[H]
\centering
\includegraphics[width=0.95\textwidth]{keygen.png}
\caption{KeyGen 命令运行截图}
\label{fig:screenshot-keygen-body}
\end{figure}

\begin{figure}[H]
\centering
\includegraphics[width=0.95\textwidth]{encrypt.png}
\caption{Encrypt 命令运行截图}
\label{fig:screenshot-encrypt-body}
\end{figure}

\begin{figure}[H]
\centering
\includegraphics[width=0.95\textwidth]{decrypt.png}
\caption{Decrypt 命令运行截图}
\label{fig:screenshot-decrypt-body}
\end{figure}

这些输出行不仅是程序提示，也对应了后续结果分析中的每个检查点。场景一、二证明满足策略时能够恢复明文；场景三、七证明不满足策略时会被拒绝；场景五证明密文完整性检查生效；场景八和九则说明新增的边界行为已经被纳入自动化测试。

真实 pairing 输出的意义在于：同一访问策略和测试集合不再只通过指数抽象模型验证，而是已经在 BLS12-381 的 $G_1/G_2/G_T$ 群元素、真实配对、哈希到群和 $G_T$ 密钥派生路径下完成闭环。由于 pairing 版密钥和密文保存真实群元素编码，输出结果同时说明 Python 后端与 PowerShell CLI 的参数转发、JSON 文件读写和错误传播能够协同工作。

\section{实验结果}

实验结果如图 \ref{fig:experiment-results} 所示。九个检查点覆盖基础策略正反例、特殊属性名、密文完整性、一般阈值门、重复属性和非法策略输入。

\begin{figure}[H]
\centering
\includegraphics[width=0.90\textwidth]{experiment_results.png}
\caption{烟雾测试场景与结果}
\label{fig:experiment-results}
\end{figure}

\subsection[场景一：属性集合 AB]{场景一：属性集合 $\{A,B\}$}

当用户拥有 A 和 B 两个属性时，用户满足策略中的 AND 分支。矩阵中 A、B 两行分别为：

\[
    M_A=(1,1),\quad M_B=(1,2).
\]

求解重构系数可得：

\[
    2M_A-M_B=2(1,1)-(1,2)=(1,0).
\]

因此 $\{A,B\}$ 可以重构秘密 $s$，进而恢复消息指数 $m$，最终解出明文。实验输出文件 \verb|out_A_B.txt| 与原始 \verb|message.txt| 完全一致。

图 \ref{fig:screenshot-out-ab-body} 展示了使用属性集合 $\{A,B\}$ 解密后的输出文件。其内容与图 \ref{fig:screenshot-message-body} 中的原始明文一致。

\begin{figure}[H]
\centering
\includegraphics[width=0.72\textwidth]{tests_artifacts_out_A_B.png}
\caption{使用属性集合 \{A,B\} 解密后的输出截图}
\label{fig:screenshot-out-ab-body}
\end{figure}

这个场景验证的是最典型的 AND 分支。由于 A、B 两行都必须参与重构，缺少任意一行都不能得到 $(1,0)$。因此它不仅证明“满足策略可以解密”，也证明系统确实在使用 LSSS 线性组合，而不是简单检查某个单一属性。

\subsection[场景二：属性集合 C]{场景二：属性集合 $\{C\}$}

属性 C 对应矩阵行：

\[
    M_C=(1,0).
\]

该行本身已经等于重构目标 $(1,0)$，因此重构系数为 $1$，用户只凭属性 C 即可满足策略。实验输出文件 \verb|out_C.txt| 与原始明文一致。

图 \ref{fig:screenshot-out-c-body} 展示了使用属性集合 $\{C\}$ 解密后的输出文件。该截图说明 OR 分支对应的单行重构路径也能正确恢复明文。

\begin{figure}[H]
\centering
\includegraphics[width=0.72\textwidth]{tests_artifacts_out_C.png}
\caption{使用属性集合 \{C\} 解密后的输出截图}
\label{fig:screenshot-out-c-body}
\end{figure}

这个场景验证的是 OR 分支。与 $\{A,B\}$ 需要两行组合不同，属性 C 对应的行本身就可以恢复秘密。这说明同一个策略中可以同时存在“多属性共同满足”和“单属性直接满足”两种路径，系统能够由矩阵自然处理二者。

\subsection[场景三：属性集合 A]{场景三：属性集合 $\{A\}$}

如果用户只有属性 A，则可用矩阵只有：

\[
    M_A=(1,1).
\]

此时需要找到 $\omega$，使得：

\[
    \omega(1,1)=(1,0).
\]

该方程要求 $\omega=1$ 且 $\omega=0$ 同时成立，显然矛盾。因此线性方程无解，系统抛出：

\begin{lstlisting}[caption={不满足策略时的错误信息}]
Attributes do not satisfy policy.
\end{lstlisting}

图 \ref{fig:screenshot-sk-a-failed-body} 展示了使用 \verb|sk_A.json| 解密时的真实终端输出。截图中的核心错误信息为 ``Attributes do not satisfy policy.''，说明系统没有为不满足访问策略的属性集合输出明文。

\begin{figure}[H]
\centering
\includegraphics[width=0.95\textwidth]{sk_A_decrypt_failed.png}
\caption{使用 sk\_A.json 解密失败的终端截图}
\label{fig:screenshot-sk-a-failed-body}
\end{figure}

这说明系统不仅能验证成功路径，也能拒绝不满足策略的私钥。

失败场景在访问控制实验中非常重要。如果只有成功解密测试，无法排除系统“过度授权”的问题。$\{A\}$ 失败说明系统没有把 AND 分支错误地当作 OR 分支处理，也说明高斯消元能够识别无解方程组。

\subsection{场景四：特殊属性名 \texttt{role:admin}}

真实系统中的属性名往往不是单个字母，而可能是 \verb|role:admin|、\verb|dept:research| 等包含冒号的字符串。Windows 文件系统中冒号具有特殊含义，如果直接将属性名拼接为私钥文件名，可能产生非预期文件。当前 CLI 会先将属性名安全化，因此 \verb|role:admin| 对应输出文件为 \verb|sk_role_admin.json|，而私钥 JSON 内部仍保留原始属性名 \verb|role:admin|。测试结果表明，该属性既能安全生成文件，也能完成解密。

测试脚本中的演示代码如下。它首先构造单属性策略 \verb|role:admin|，随后检查私钥文件是否被安全命名为 \verb|sk_role_admin.json|，最后使用该私钥完成加解密。

\begin{lstlisting}[language=bash,caption={特殊属性名场景测试代码}]
[pscustomobject]@{ attr = "role:admin" } |
  ConvertTo-Json -Depth 10 |
  Set-Content -LiteralPath $specialPolicyPath -Encoding UTF8

& $bin keygen -Pub (Join-Path $art "public.json") `
  -Msk (Join-Path $art "master.json") `
  -OutDir $art -Attrs "role:admin"

$specialSkPath = Join-Path $art "sk_role_admin.json"
if (-not (Test-Path -LiteralPath $specialSkPath)) {
    throw "Scenario 4 failed: filename-safe secret key was not created."
}

& $bin encrypt -Pub (Join-Path $art "public.json") `
  -Policy $specialPolicyPath -In (Join-Path $root "examples\message.txt") `
  -Out $specialCtPath
& $bin decrypt -Pub (Join-Path $art "public.json") `
  -Sk $specialSkPath -In $specialCtPath -Out $specialOutPath
\end{lstlisting}

这个场景属于工程健壮性测试。真实属性名通常不会只由单个大写字母组成，可能包含部门、角色、等级等命名空间信息。文件名安全化如果处理不当，可能导致私钥文件无法生成，或者生成的文件名和属性语义不一致。该测试确认安全化只影响文件名，不影响策略匹配。

\subsection{场景五：密文篡改检测}

为验证对称载荷的完整性保护，测试脚本将已生成密文中的 \verb|payload.tag| 替换为错误 HMAC 标签，再尝试解密。系统在 AES 解密前重新计算 HMAC，并抛出：

\begin{lstlisting}[caption={密文篡改时的错误信息}]
Invalid ciphertext tag.
\end{lstlisting}

测试脚本中的核心做法是读取原始密文，替换认证标签字段，再尝试用满足策略的 A、B 私钥解密。由于 HMAC 校验失败，程序必须拒绝输出明文。

\begin{lstlisting}[language=bash,caption={密文篡改场景测试代码}]
$ctObj = Get-Content -Raw -LiteralPath (Join-Path $art "ct.json") `
  -Encoding UTF8 | ConvertFrom-Json
$ctObj.payload.tag = [Convert]::ToBase64String((New-Object byte[] 32))
$ctObj | ConvertTo-Json -Depth 100 |
  Set-Content -LiteralPath $tamperedCtPath -Encoding UTF8

try {
    & $bin decrypt -Pub (Join-Path $art "public.json") `
      -Sk (Join-Path $art "sk_A_B.json") `
      -In $tamperedCtPath -Out (Join-Path $art "out_tampered.txt")
    throw "Scenario 5 failed: tampered ciphertext should NOT decrypt."
} catch {
    if ($_.Exception.Message -notmatch "Invalid ciphertext tag") {
        throw "Scenario 5 failed: unexpected error."
    }
}
\end{lstlisting}

这说明密文载荷被篡改时不会被静默接受。

该测试也说明 CP-ABE 层和对称加密层之间的边界是清楚的。即使用户属性满足策略，如果密文载荷的认证标签错误，系统仍会拒绝输出明文。这避免了“访问策略满足就忽略密文完整性”的错误实现。

\subsection{场景六与七：2-of-3 阈值门}

为验证访问控制树中一般阈值门的处理能力，测试脚本构造了 $2$-of-$3(D,E,F)$ 策略。属性集合 $\{D,F\}$ 满足该策略，因此可以解密；属性集合 $\{D\}$ 只包含一个属性，不满足 $2$-of-$3$ 条件，因此被拒绝。这一测试表明，系统并非只支持 AND/OR 特例，也能处理一般阈值门。

测试脚本中的关键代码如下。前半部分构造阈值策略并验证 $\{D,F\}$ 成功；后半部分只为属性 D 生成私钥，并要求系统拒绝解密。

\begin{lstlisting}[language=bash,caption={2-of-3 阈值门成功与失败测试代码}]
[pscustomobject]@{
    k = 2
    children = @(
        [pscustomobject]@{ attr = "D" },
        [pscustomobject]@{ attr = "E" },
        [pscustomobject]@{ attr = "F" }
    )
} | ConvertTo-Json -Depth 10 |
  Set-Content -LiteralPath $thresholdPolicyPath -Encoding UTF8

& $bin keygen -Pub (Join-Path $art "public.json") `
  -Msk (Join-Path $art "master.json") -OutDir $art -Attrs "D,F"
& $bin encrypt -Pub (Join-Path $art "public.json") `
  -Policy $thresholdPolicyPath -In (Join-Path $root "examples\message.txt") `
  -Out $thresholdCtPath
& $bin decrypt -Pub (Join-Path $art "public.json") `
  -Sk (Join-Path $art "sk_D_F.json") `
  -In $thresholdCtPath -Out $thresholdOutPath

& $bin keygen -Pub (Join-Path $art "public.json") `
  -Msk (Join-Path $art "master.json") -OutDir $art -Attrs "D"
try {
    & $bin decrypt -Pub (Join-Path $art "public.json") `
      -Sk (Join-Path $art "sk_D.json") `
      -In $thresholdCtPath -Out (Join-Path $art "out_D.txt")
    throw "Scenario 7 failed: {D} should NOT satisfy 2-of-3 policy."
} catch {
    if ($_.Exception.Message -notmatch "Attributes do not satisfy") {
        throw "Scenario 7 failed: unexpected error."
    }
}
\end{lstlisting}

这两个场景扩展了示例策略的覆盖范围。AND 和 OR 都可以看成阈值门的特例，但如果不单独测试 $2$-of-$3$，仍难以说明实现已经覆盖一般阈值门路径。通过同时测试 $\{D,F\}$ 成功和 $\{D\}$ 失败，可以确认阈值条件不是被简单地写死为“至少一个”或“全部”。

\subsection{场景八：重复属性}

重复属性测试构造策略 $X\wedge X$，其 LSSS 矩阵中会出现多行映射到同一属性 X。解密端筛选行时会把所有 $\rho(i)=X$ 的行都纳入线性方程求解，并复用私钥中的同一个 $K_X$ 分量完成计算。该测试用于验证系统能够处理“多行映射到同一属性”的工程场景。

测试脚本中的核心代码如下。它先构造两个叶子均为 X 的 $2$-of-$2$ 策略，再只为属性 X 生成私钥；如果系统没有正确处理重复属性行，该场景会在解密或明文比对阶段失败。

\begin{lstlisting}[language=bash,caption={重复属性场景测试代码}]
[pscustomobject]@{
    k = 2
    children = @(
        [pscustomobject]@{ attr = "X" },
        [pscustomobject]@{ attr = "X" }
    )
} | ConvertTo-Json -Depth 10 | Set-Content -LiteralPath $duplicatePolicyPath -Encoding UTF8

& $bin keygen -Pub (Join-Path $art "public.json") -Msk (Join-Path $art "master.json") -OutDir $art -Attrs "X"
& $bin encrypt -Pub (Join-Path $art "public.json") -Policy $duplicatePolicyPath -In (Join-Path $root "examples\message.txt") -Out $duplicateCtPath
& $bin decrypt -Pub (Join-Path $art "public.json") -Sk (Join-Path $art "sk_X.json") -In $duplicateCtPath -Out $duplicateOutPath

$outDuplicate = Get-Content -Raw -LiteralPath $duplicateOutPath -Encoding UTF8
if ($orig -ne $outDuplicate) {
    throw "Scenario 8 failed: repeated attribute rows should decrypt with attribute X."
}
\end{lstlisting}

\subsection{场景九：非法策略输入}

非法策略测试包含多类错误访问树：阈值 $k=0$、$k$ 大于子节点数量、缺失 $k$、子节点列表为空、属性名为空，以及叶子节点同时携带 \verb|children|。系统在树到 LSSS 矩阵转换前执行递归校验，并在发现非法结构时抛出 \verb|Invalid policy| 相关错误。这样可以避免非法访问树进入矩阵构造流程，减少后续难以定位的运行时异常。

非法策略测试不是只检查单个错误输入，而是把多类错误访问树组织成数组逐个加密。实验现象是每个非法策略都在加密阶段抛出 \verb|Invalid policy| 相关错误，且没有任何非法策略生成可用密文。

\begin{lstlisting}[language=bash,caption={非法策略输入场景测试代码}]
$invalidPolicies = @(
    [pscustomobject]@{ Name = "k0"; Policy = [pscustomobject]@{ k = 0; children = @([pscustomobject]@{ attr = "Z" }) } },
    [pscustomobject]@{ Name = "k_too_large"; Policy = [pscustomobject]@{ k = 3; children = @([pscustomobject]@{ attr = "Y" }, [pscustomobject]@{ attr = "Z" }) } },
    [pscustomobject]@{ Name = "missing_k"; Policy = [pscustomobject]@{ children = @([pscustomobject]@{ attr = "Z" }) } },
    [pscustomobject]@{ Name = "empty_children"; Policy = [pscustomobject]@{ k = 1; children = @() } },
    [pscustomobject]@{ Name = "empty_attr"; Policy = [pscustomobject]@{ attr = "" } },
    [pscustomobject]@{ Name = "leaf_with_children"; Policy = [pscustomobject]@{ attr = "Q"; children = @([pscustomobject]@{ attr = "R" }) } }
)

$invalidPoliciesRejected = 0
foreach ($case in $invalidPolicies) {
    $invalidPolicyPath = Join-Path $art "policy_invalid_$($case.Name).json"
    $case.Policy | ConvertTo-Json -Depth 10 | Set-Content -LiteralPath $invalidPolicyPath -Encoding UTF8
    try {
        & $bin encrypt -Pub (Join-Path $art "public.json") -Policy $invalidPolicyPath -In (Join-Path $root "examples\message.txt") -Out (Join-Path $art "ct_invalid.json")
        throw "Invalid policy should not encrypt: $invalidPolicyPath"
    } catch {
        if ($_.Exception.Message -notmatch "Invalid policy") {
            throw "Scenario 9 failed: unexpected error for invalid policy."
        }
        $invalidPoliciesRejected++
    }
}
\end{lstlisting}

通过把非法策略统一放入数组，测试脚本保留了继续扩展的空间，例如非整数 \verb|k|、空根节点、缺失 \verb|children| 或更复杂的畸形 JSON。当前实现已经覆盖了最容易导致矩阵构造错误的几类输入，因此能更早暴露策略格式问题。

\section{结果分析}

从实验现象看，九个场景形成了一条较完整的证据链。场景一和场景二分别对应示例策略 $(A\wedge B)\vee C$ 的 AND 分支和 OR 分支，二者都能恢复原始明文，说明系统生成的密钥分量、密文分量和对称载荷可以完成一次完整的加密解密闭环。场景三使用属性集合 $\{A\}$ 尝试解密并被拒绝，说明系统没有把 ``拥有部分属性'' 错误地视为满足策略。成功路径和失败路径同时成立，才能说明访问控制不是形式上的字符串比较，而是由矩阵重构条件决定。

从 LSSS 角度看，实验结果说明访问树语义与线性重构语义保持一致。加密阶段把访问树转换为矩阵 $(M,\rho)$，解密阶段根据用户属性集合 $S$ 选取行集合 $I=\{i:\rho(i)\in S\}$，再求解 $M_I^T\omega=e_1$。当属性集合为 $\{A,B\}$ 时，A 行和 B 行可以共同重构第一维；当属性集合为 $\{C\}$ 时，C 行单独即可完成重构；当属性集合为 $\{A\}$ 时，可用行不足以重构第一维，因此线性方程无解。由此可以看出，程序判断策略满足性时依赖的是 LSSS 的线性代数条件，而不是针对示例策略手写 AND/OR 判断分支。

一般阈值门测试进一步扩展了这一结论。$2$-of-$3(D,E,F)$ 场景中，$\{D,F\}$ 能够解密而 $\{D\}$ 被拒绝，说明系统处理的不是简单的 ``至少一个属性满足'' 或 ``全部属性满足'' 两种特例，而是按照阈值 $k$ 和子节点数 $n$ 共同决定访问条件。该结果与第三章中访问树一般性的设计相互对应：内部节点统一保存 $k$ 与 \verb|children|，矩阵构造阶段为阈值门引入 $k-1$ 个随机维度，解密阶段再由行集合可重构性给出最终判定。

工程边界测试说明实验对象不仅停留在理想公式输入上。特殊属性名 \verb|role:admin| 能够生成安全文件名并保留原始属性语义，说明 CLI 文件系统处理没有破坏策略匹配；密文篡改测试在 HMAC 检查阶段被拒绝，说明满足访问策略并不会绕过对称载荷完整性验证；重复属性 $X\wedge X$ 能够成功重构，说明 $\rho$ 允许多行映射到同一属性，解密端没有错误地只保留某个属性的第一行；非法策略输入在加密前被拒绝，说明策略校验可以在矩阵构造前截断明显错误的访问树结构。

从实验覆盖关系看，各场景可以归纳为四类。第一类是正确性场景，验证满足策略时明文能够被恢复；第二类是拒绝场景，验证不满足策略或密文被篡改时不会输出明文；第三类是表达能力场景，验证一般阈值门和重复属性行能够被 LSSS 机制处理；第四类是输入健壮性场景，验证特殊属性名和非法策略不会破坏系统运行。四类结果共同支撑本章结论：本项目实现的核心并非一次孤立的文件加密演示，而是围绕访问控制树、LSSS 矩阵、属性私钥和线性重构建立了可复现实验链路。

结合课程题目要求，这组实验能够说明系统已经完成两个主要目标：其一，访问控制树能够自动转换为 LSSS 矩阵，并在密文中以可检查形式保存；其二，Setup、KeyGen、Encrypt 和 Decrypt 四个算法能够围绕该矩阵形成完整流程，并区分满足策略与不满足策略的属性集合。与实验三课件~[1]中关于 Waters 2011 方案~[3]和 FABEO 方案~[5]的加解密测试要求相比，本项目的检验重点更集中在 LSSS-CP-ABE 的访问结构语义上。也就是说，本实验不只比较一次明文和解密输出是否相同，而是通过多种属性集合和边界场景验证矩阵重构条件与访问树语义的一致性。

需要说明的是，本章实验结论分为两个层次。指数抽象版验证了 LSSS-CP-ABE 的矩阵构造、份额重构和访问控制行为，便于解释公式中的代数抵消；真实 pairing 版进一步验证了同一访问结构可以在 BLS12-381 群元素和真实配对运算上运行，说明底层群运算真实化后核心链路仍然成立。二者仍不能被直接解释为完整安全证明，因为形式化安全还涉及所选论文模型、随机数采样、哈希到群、序列化检查、参数选择和攻击模型等条件。复杂嵌套阈值树、更多畸形 JSON 结构和更严格的密文结构校验仍属于后续可继续扩展的验证对象。
